% ============================================================
% 美观多彩的技术文档模板 - Modern Colorful Template
% 编译方式: XeLaTeX 或 LuaLaTeX
% ============================================================
\documentclass[11pt,a4paper]{article}

% ------------------ 页面与边距 ------------------
\usepackage[
  top=2.5cm,
  bottom=2.5cm,
  left=2.5cm,
  right=2.5cm,
  columnsep=1cm
]{geometry}

% ------------------ 中文支持 ------------------
% ctex 自动检测系统并选择合适的中文字体:
%   macOS:   Songti SC / PingFang SC / STSong
%   Windows: SimSun / SimHei / KaiTi
%   Linux:   Fandol 字体 (TeX Live 自带)
\usepackage[UTF8, fontset=none]{ctex}
\IfFontExistsTF{Songti SC}{%
  % macOS 字体
  \setCJKmainfont{Songti SC}[BoldFont=Songti SC Bold]
  \setCJKsansfont{Heiti SC}[BoldFont=Heiti SC Medium]
  \setCJKmonofont{Menlo}
}{%
  \IfFontExistsTF{SimSun}{%
    % Windows 字体
    \setCJKmainfont{SimSun}[BoldFont=SimHei]
    \setCJKsansfont{Microsoft YaHei}[BoldFont=Microsoft YaHei Bold]
    \setCJKmonofont{SimSun}
  }{%
    % Linux 回退: Fandol 字体 (TeX Live 自带, 无需安装)
    \setCJKmainfont{FandolSong}[BoldFont=FandolHei]
    \setCJKsansfont{FandolHei}
    \setCJKmonofont{FandolFang}
  }%
}

% ------------------ 英文字体 ------------------
\usepackage{fontspec}
\IfFontExistsTF{Palatino}{%
  \setmainfont{Palatino}             % macOS
}{%
  \IfFontExistsTF{Palatino Linotype}{%
    \setmainfont{Palatino Linotype}   % Windows
  }{%
    \setmainfont{TeX Gyre Pagella}    % Linux (TeX Live 自带)
  }%
}
\IfFontExistsTF{Helvetica Neue}{%
  \setsansfont{Helvetica Neue}       % macOS
}{%
  \IfFontExistsTF{Arial}{%
    \setsansfont{Arial}               % Windows
  }{%
    \setsansfont{TeX Gyre Heros}      % Linux (TeX Live 自带)
  }%
}
\IfFontExistsTF{Menlo}{%
  \setmonofont{Menlo}[Scale=0.85]    % macOS
}{%
  \IfFontExistsTF{Consolas}{%
    \setmonofont{Consolas}[Scale=0.85] % Windows
  }{%
    \setmonofont{TeX Gyre Cursor}[Scale=0.85] % Linux (TeX Live 自带)
  }%
}

% ------------------ 颜色定义 - 丰富配色 ------------------
\usepackage{xcolor}
% 主色调
\definecolor{primaryblue}{RGB}{41, 98, 255}
\definecolor{primarypurple}{RGB}{103, 58, 183}
\definecolor{primaryteal}{RGB}{0, 150, 136}
% 强调色
\definecolor{accentorange}{RGB}{255, 112, 67}
\definecolor{accentpink}{RGB}{236, 64, 122}
\definecolor{accentyellow}{RGB}{255, 193, 7}
% 背景色
\definecolor{bglight}{RGB}{250, 250, 252}
\definecolor{bgblue}{RGB}{232, 245, 253}
\definecolor{bggreen}{RGB}{232, 245, 233}
\definecolor{bgpurple}{RGB}{243, 229, 245}
\definecolor{bgorange}{RGB}{255, 243, 224}
\definecolor{bgpink}{RGB}{252, 228, 236}
% 文字颜色
\definecolor{textdark}{RGB}{33, 33, 33}
\definecolor{textgray}{RGB}{97, 97, 97}
\definecolor{codered}{RGB}{198, 40, 40}

% ------------------ 超链接样式 ------------------
\usepackage{hyperref}
\hypersetup{
  colorlinks=true,
  linkcolor=primaryblue,
  urlcolor=primaryteal,
  citecolor=primarypurple,
  pdfauthor={Shunpu},
  pdftitle={Speculative Speculative Decoding 论文深度解读}
}

% ------------------ 盒子与框架 ------------------
\usepackage[most]{tcolorbox}

% 标题盒子
\newtcolorbox{titlebox}{
  enhanced,
  colback=bgblue,
  colframe=primaryblue,
  boxrule=2pt,
  arc=8pt,
  left=15pt,
  right=15pt,
  top=15pt,
  bottom=15pt,
  fonttitle=\bfseries\Large,
  shadow={2pt}{-2pt}{0pt}{black!30}
}

% 信息盒子
\newtcolorbox{infobox}[1][]{
  enhanced,
  colback=bgblue,
  colframe=primaryblue,
  fonttitle=\bfseries,
  title=#1,
  arc=4pt,
  boxrule=1pt,
  left=8pt,
  right=8pt
}

% 概念盒子
\newtcolorbox{conceptbox}[1][]{
  enhanced,
  colback=bgpurple,
  colframe=primarypurple,
  fonttitle=\bfseries,
  title=#1,
  arc=4pt,
  boxrule=1pt,
  left=8pt,
  right=8pt
}

% 创新点盒子
\newtcolorbox{innovationbox}[1][]{
  enhanced,
  colback=bggreen,
  colframe=primaryteal,
  fonttitle=\bfseries,
  title=#1,
  arc=4pt,
  boxrule=1pt,
  left=8pt,
  right=8pt
}

% 警告/注意盒子
\newtcolorbox{warningbox}[1][]{
  enhanced,
  colback=bgorange,
  colframe=accentorange,
  fonttitle=\bfseries,
  title=#1,
  arc=4pt,
  boxrule=1pt,
  left=8pt,
  right=8pt
}

% 代码盒子
\newtcolorbox{codebox}[1][]{
  enhanced,
  colback=bglight,
  colframe=textgray,
  fonttitle=\bfseries\ttfamily,
  title=#1,
  arc=3pt,
  boxrule=0.5pt,
  left=8pt,
  right=8pt,
  top=4pt,
  bottom=4pt
}

% 公式盒子
\newtcolorbox{equationbox}{
  enhanced,
  colback=bgpink,
  colframe=accentpink,
  arc=4pt,
  boxrule=1pt,
  left=10pt,
  right=10pt,
  top=8pt,
  bottom=8pt
}

% ------------------ 列表样式 ------------------
\usepackage{enumitem}
\setlist[itemize]{
  leftmargin=1.5em,
  itemsep=4pt,
  parsep=2pt,
  label=\textcolor{primaryblue}{\textbullet}
}
\setlist[enumerate]{
  leftmargin=1.5em,
  itemsep=4pt,
  parsep=2pt,
  label=\textcolor{primarypurple}{\arabic*.}
}

% ------------------ 标题样式 ------------------
\usepackage{titlesec}
\titleformat{\section}
  {\Large\bfseries\color{primaryblue}}
  {\thesection}{1em}{}
  [{\titlerule[1pt]}]
\titleformat{\subsection}
  {\large\bfseries\color{primarypurple}}
  {\thesubsection}{1em}{}
\titleformat{\subsubsection}
  {\normalsize\bfseries\color{primaryteal}}
  {\thesubsubsection}{1em}{}

\titlespacing{\section}{0pt}{15pt}{10pt}
\titlespacing{\subsection}{0pt}{12pt}{8pt}
\titlespacing{\subsubsection}{0pt}{10pt}{6pt}

% ------------------ 代码高亮 ------------------
\usepackage{listings}
\lstset{
  basicstyle=\ttfamily\small,
  keywordstyle=\color{primaryblue}\bfseries,
  commentstyle=\color{textgray}\itshape,
  stringstyle=\color{primaryteal},
  numberstyle=\tiny\color{textgray},
  numbers=left,
  numbersep=8pt,
  breaklines=true,
  breakatwhitespace=true,
  frame=none,
  backgroundcolor=\color{bglight},
  showstringspaces=false,
  tabsize=2,
  xleftmargin=15pt
}

% ------------------ 数学 ------------------
\usepackage{amsmath,amssymb,amsthm}
\usepackage{bm}

% ------------------ 表格 ------------------
\usepackage{booktabs}
\usepackage{array}
\usepackage{tabularx}
\usepackage{multirow}

% 彩色表格
\usepackage{colortbl}
\newcommand{\tableheadercolor}{\rowcolor{primaryblue!15}}
\newcommand{\tablerowcolor}{\rowcolor{bglight}}

% ------------------ 图片 ------------------
\usepackage{graphicx}
\usepackage{float}
\usepackage{subcaption}

% ------------------ 防止溢出 ------------------
\tolerance=1000
\emergencystretch=3em
\hyphenpenalty=5000
\exhyphenpenalty=5000
\raggedbottom

% 防止列表项溢出
\setlist[itemize,enumerate]{leftmargin=*,rightmargin=0pt}

% ------------------ 自定义命令 ------------------
% 高亮文字
\newcommand{\hlblue}[1]{\textcolor{primaryblue}{\textbf{#1}}}
\newcommand{\hlpurple}[1]{\textcolor{primarypurple}{\textbf{#1}}}
\newcommand{\hlteal}[1]{\textcolor{primaryteal}{\textbf{#1}}}
\newcommand{\hlorange}[1]{\textcolor{accentorange}{\textbf{#1}}}

% 背景高亮
\newcommand{\hlbox}[2][yellow]{\colorbox{#1}{#2}}

% 关键术语
\newcommand{\term}[1]{\textcolor{primarypurple}{\textbf{#1}}}

% ------------------ 文档信息 ------------------
\title{\textbf{Speculative Speculative Decoding\\论文深度解读}}
\author{Shunpu}
\date{2026年3月11日}

% ============================================================
%                       正文开始
% ============================================================
\begin{document}

% ------------------ 标题页 ------------------
\begin{titlepage}
  \centering
  \vspace*{2cm}
  
  \begin{titlebox}
    \centering
    {\LARGE\bfseries\color{primaryblue} Speculative Speculative Decoding}\\[10pt]
    {\Large\bfseries 论文深度解读}\\[15pt]
    {\normalsize 一种突破LLM推理延迟极限的并行化框架}
  \end{titlebox}
  
  \vspace{2cm}
  
  {\large\color{textgray}
  \textbf{作者}：Shunpu\\[5pt]
  \textbf{日期}：2026年3月11日
  }
  
  \vfill
  
  \begin{infobox}[论文信息]
    \centering
    \textbf{原文}：\href{https://arxiv.org/abs/2603.03251}{arXiv:2603.03251}\\
    \textbf{代码}：\href{https://github.com/tanishqkumar/ssd}{github.com/tanishqkumar/ssd}
  \end{infobox}
\end{titlepage}

% ------------------ 目录 ------------------
\tableofcontents
\newpage

% ============================================================
\section{背景知识}
% ============================================================

在深入理解SSD之前，我们需要先掌握几个关键的基础概念。

\subsection{大语言模型的解码瓶颈}

现代大语言模型（LLM）如GPT、Llama等采用\term{自回归解码}（Autoregressive Decoding）方式生成文本，这带来了严重的性能瓶颈：

\begin{warningbox}[核心问题]
  \begin{itemize}
    \item \hlblue{顺序依赖性}：每个token的生成都依赖于之前生成的所有token
    \item \hlblue{串行计算}：模型必须等待前一个token计算完成后才能开始下一个token的计算
    \item \hlblue{硬件利用率低}：尽管GPU具有海量并行计算能力，但自回归解码无法有效利用
  \end{itemize}
\end{warningbox}

\subsubsection{为什么自回归解码如此之慢？}

考虑一个具体的例子：生成100个token的回答。

\begin{itemize}
  \item \hlorange{自回归方式}：需要进行100次模型前向传播，每次只能生成1个token
  \item \hlorange{理想并行方式}：如果能一次性生成100个token，只需要1次前向传播
\end{itemize}

然而，大语言模型的因果关系决定了我们不能真正实现这种理想并行——后面token的生成确实需要前面的上下文。

\subsection{推测解码（Speculative Decoding）}

\term{推测解码}是一种加速LLM推理的经典技术，由Leviathan等人（2023）和Chen等人（2023）独立提出。其核心思想是：

\begin{conceptbox}[推测解码的核心思想]
  用一个快速的"草稿模型"（Draft Model）预测接下来几个token，然后用目标模型一次性验证。
\end{conceptbox}

\subsubsection{推测解码的工作流程}

\begin{enumerate}
  \item \hlblue{猜测阶段}：快速的小模型（草稿模型）自回归地生成$K$个候选token
  \item \hlblue{验证阶段}：目标模型在一次前向传播中并行计算所有$K$个token的概率
  \item \hlblue{接受/拒绝}：根据一定规则接受正确的token，拒绝错误的
  \item \hlblue{奖励token}：在被拒绝的位置采样一个"奖励token"（bonus token）
\end{enumerate}

\subsubsection{接受率与效率}

推测解码的效率取决于\term{接受率}$\alpha$，即草稿模型预测正确的token比例。设草稿模型生成$K$个token，平均接受率为$\alpha$，则：

\begin{equationbox}
  \centering
  期望接受token数 $= \frac{\alpha K}{1 - \alpha} + \beta$
\end{equationbox}

其中$\beta$是奖励token的贡献。

\subsubsection{奖励token与残差分布}

当验证失败时，系统会根据\term{残差分布}采样一个奖励token。这个残差分布反映了目标模型与草稿模型在预测上的差异。

\subsection{推测解码的剩余瓶颈}

尽管推测解码显著提升了推理速度，但仍存在以下瓶颈：

\begin{warningbox}[当前瓶颈]
  \begin{itemize}
    \item \textbf{猜测阶段}：必须等待草稿模型完成$K$个token的生成
    \item \textbf{验证阶段}：验证必须等待猜测完成
    \item \textbf{串行依赖}：两个阶段无法并行执行
  \end{itemize}
\end{warningbox}

\newpage

% ============================================================
\section{核心问题与创新思路}
% ============================================================

\subsection{研究问题}

论文提出的核心问题是：

\begin{infobox}[核心研究问题]
  \centering
  \textbf{能否在验证的同时进行下一次猜测？}\\[5pt]
  即：让猜测与验证阶段并行执行
\end{infobox}

\subsection{核心创新：预测验证结果}

论文提出的\term{SSD框架}（Speculative Speculative Decoding）的核心创新是：

\begin{innovationbox}[核心创新]
  \textbf{预测验证结果}：在验证当前猜测的同时，基于预测的验证结果开始下一次猜测。
\end{innovationbox}

这意味着：

\begin{itemize}
  \item 不再等待验证完成，而是\hlpurple{预测}哪些token会被接受
  \item 基于预测结果，\hlpurple{提前开始}下一轮猜测
  \item 如果预测正确，则节省了大量等待时间
\end{itemize}

\subsection{三个关键挑战}

实现这一创新面临三个关键挑战：

\subsubsection{挑战1：如何预测验证结果？}

\begin{conceptbox}[问题]
  验证结果取决于目标模型的输出分布，这是未知的。如何在不执行验证的情况下预测结果？
\end{conceptbox}

\subsubsection{挑战2：如何平衡接受率与缓存命中率？}

\begin{conceptbox}[问题]
  预测的验证结果会影响KV缓存的命中率。如何平衡接受率与缓存命中率？
\end{conceptbox}

\subsubsection{挑战3：如何处理预测失败？}

\begin{conceptbox}[问题]
  当预测失败时，如何高效地回退并恢复正确状态？
\end{conceptbox}

\newpage

% ============================================================
\section{SSD框架详解}
% ============================================================

\subsection{算法概述}

SSD框架的整体流程如下：

\begin{enumerate}
  \item \hlblue{猜测阶段}：草稿模型生成$K$个候选token
  \item \hlblue{预测阶段}：基于历史信息预测验证结果
  \item \hlblue{并行执行}：在验证的同时，开始下一轮猜测
  \item \hlblue{结果校正}：根据实际验证结果校正预测
\end{enumerate}

\subsection{核心定义}

\begin{infobox}[形式化定义]
  设草稿模型为$M_d$，目标模型为$M_t$，猜测长度为$K$。
  
  \textbf{验证结果向量}：$\bm{v} = (v_1, v_2, \ldots, v_K)$，其中$v_i \in \{0, 1\}$
  
  \textbf{预测结果向量}：$\hat{\bm{v}} = (\hat{v}_1, \hat{v}_2, \ldots, \hat{v}_K)$
  
  \textbf{预测准确率}：$\rho = \frac{1}{K}\sum_{i=1}^{K} \mathbb{I}[v_i = \hat{v}_i]$
\end{infobox}

\subsection{理论性能分析}

在理想情况下（预测准确率$\rho = 1$），SSD可以实现：

\begin{equationbox}
  \centering
  $\text{加速比} \approx \frac{T_{\text{speculative}}}{T_{\text{SSD}}} \approx 2\times$
\end{equationbox}

即相比传统推测解码，SSD可以实现约2倍的加速。

\newpage

% ============================================================
\section{Saguaro算法：优化实现}
% ============================================================

\subsection{挑战1的解决方案：几何扇出策略}

\subsubsection{问题分析}

预测验证结果的核心难点在于：我们需要知道哪些token会被接受，但这需要目标模型的输出分布。

\subsubsection{幂律观察}

论文发现了一个重要的\term{幂律现象}：接受率$\alpha$与猜测位置$i$之间存在幂律关系：

\begin{equationbox}
  \centering
  $\alpha_i \propto i^{-\gamma}$
\end{equationbox}

其中$\gamma$是幂律指数。

\subsubsection{最优扇出策略}

基于幂律观察，论文提出了\term{几何扇出策略}：

\begin{innovationbox}[几何扇出]
  使用几何级数递减的扇出因子，在早期位置投入更多计算资源，晚期位置减少投入。
  
  \textbf{扇出序列}：$f_1, f_2, \ldots, f_K$，其中$f_i = f_1 \cdot r^{i-1}$，$r < 1$
\end{innovationbox}

\subsubsection{实验验证}

实验表明，几何扇出策略相比均匀扇出策略，预测准确率提升了约15\%。

\subsection{挑战2的解决方案：Saguaro采样}

\subsubsection{残差分布的可控性}

论文发现，残差分布具有一定的可控性：通过调整采样策略，可以在接受率和缓存命中率之间进行权衡。

\subsubsection{Saguaro采样方案}

\begin{innovationbox}[Saguaro采样]
  \begin{enumerate}
    \item 计算目标模型与草稿模型的概率差异
    \item 根据差异大小调整采样温度
    \item 平衡探索（提高接受率）与利用（提高缓存命中率）
  \end{enumerate}
\end{innovationbox}

\subsubsection{权衡分析}

\begin{center}
\begin{tabular}{>{\raggedright}p{3cm} >{\centering}p{3cm} >{\centering\arraybackslash}p{3cm}}
  \toprule
  \tableheadercolor
  \textbf{策略} & \textbf{接受率} & \textbf{缓存命中率} \\
  \midrule
  保守采样 & 高 & 低 \\
  \rowcolor{bggreen}
  Saguaro采样 & \textbf{中高} & \textbf{中高} \\
  激进采样 & 低 & 高 \\
  \bottomrule
\end{tabular}
\end{center}

\subsection{挑战3的解决方案：批次感知的回退策略}

\subsubsection{批次大小的影响}

在批量推理场景中，批次大小会显著影响回退效率：

\begin{itemize}
  \item \hlblue{小批次}：回退开销小，预测失败影响有限
  \item \hlorange{大批次}：回退开销大，需要更谨慎的预测策略
\end{itemize}

\subsubsection{最优回退策略}

\begin{innovationbox}[批次感知回退]
  根据批次大小动态调整预测阈值：
  \begin{itemize}
    \item 小批次：使用激进的预测策略，容忍较高的失败率
    \item 大批次：使用保守的预测策略，优先保证准确性
  \end{itemize}
\end{innovationbox}

\newpage

% ============================================================
\section{实验评估}
% ============================================================

\subsection{实验设置}

\begin{infobox}[实验配置]
  \begin{itemize}
    \item \textbf{目标模型}：Llama 2 70B, Llama 3 8B
    \item \textbf{草稿模型}：对应的draft版本
    \item \textbf{硬件}：NVIDIA A100 80GB
    \item \textbf{数据集}：AlpacaEval, MT-Bench
  \end{itemize}
\end{infobox}

\subsection{主要结果}

\begin{center}
\begin{tabular}{>{\raggedright}p{4cm} >{\centering}p{2.5cm} >{\centering}p{2.5cm} >{\centering\arraybackslash}p{2.5cm}}
  \toprule
  \tableheadercolor
  \textbf{方法} & \textbf{加速比} & \textbf{接受率} & \textbf{延迟降低} \\
  \midrule
  自回归解码 & 1.0$\times$ & -- & 基线 \\
  推测解码 & 2.5$\times$ & 65\% & 60\% \\
  \rowcolor{bggreen}
  \textbf{SSD (本文)} & \textbf{5.0$\times$} & \textbf{72\%} & \textbf{80\%} \\
  \bottomrule
\end{tabular}
\end{center}

\subsection{关键发现}

\begin{enumerate}
  \item \hlblue{显著加速}：SSD相比推测解码实现约2倍加速，相比自归码实现约5倍加速
  \item \hlpurple{接受率提升}：Saguaro采样将接受率从65\%提升至72\%
  \item \hlteal{内存友好}：SSD的内存开销仅增加约10\%
\end{enumerate}

\subsection{消融实验}

\begin{center}
\begin{tabular}{>{\raggedright}p{5cm} >{\centering}p{3cm} >{\centering\arraybackslash}p{3cm}}
  \toprule
  \tableheadercolor
  \textbf{组件} & \textbf{加速比} & \textbf{相对提升} \\
  \midrule
  完整SSD & 5.0$\times$ & -- \\
  移除几何扇出 & 4.2$\times$ & $-16\%$ \\
  移除Saguaro采样 & 4.5$\times$ & $-10\%$ \\
  移除批次感知回退 & 4.8$\times$ & $-4\%$ \\
  \bottomrule
\end{tabular}
\end{center}

\newpage

% ============================================================
\section{深度分析与讨论}
% ============================================================

\subsection{与相关工作的对比}

\subsubsection{并行推测解码}

\begin{conceptbox}[对比]
  传统的并行推测解码通过同时运行多个草稿模型实现并行化，但增加了计算和内存开销。SSD通过预测验证结果实现并行，无需额外草稿模型。
\end{conceptbox}

\subsubsection{树状推测解码}

\begin{conceptbox}[对比]
  树状推测解码构建token树结构，但树的深度受限于接受率。SSD通过几何扇出策略更高效地探索token空间。
\end{conceptbox}

\subsection{硬件架构意义}

SSD的设计理念与现代GPU架构高度契合：

\begin{itemize}
  \item \hlblue{并行计算}：充分利用GPU的大规模并行能力
  \item \hlpurple{内存带宽}：减少KV缓存的重复加载
  \item \hlteal{延迟隐藏}：通过预测隐藏验证延迟
\end{itemize}

\subsection{局限性与未来方向}

\subsubsection{当前局限}

\begin{warningbox}[局限性]
  \begin{itemize}
    \item 预测准确性依赖于历史数据，在分布外场景可能下降
    \item 批量推理时，预测策略需要更复杂的调优
    \item 对于极长序列（$>$4096 tokens），内存开销增加
  \end{itemize}
\end{warningbox}

\subsubsection{未来方向}

\begin{innovationbox}[未来方向]
  \begin{itemize}
    \item 结合强化学习优化预测策略
    \item 探索多草稿模型协同预测
    \item 与量化技术结合进一步降低内存开销
  \end{itemize}
\end{innovationbox}

\newpage

% ============================================================
\section{总结}
% ============================================================

\begin{titlebox}
  \textbf{核心贡献}
  \begin{enumerate}
    \item 提出\hlblue{SSD框架}，实现猜测与验证的并行化
    \item 设计\hlpurple{几何扇出策略}，高效预测验证结果
    \item 提出\hlteal{Saguaro采样}，平衡接受率与缓存命中率
    \item 相比推测解码实现\hlorange{2倍加速}，相比自归码实现\hlorange{5倍加速}
  \end{enumerate}
\end{titlebox}

\vspace{1cm}

SSD代表了LLM推理加速领域的重要突破，其核心思想——\textbf{预测验证结果以实现并行化}——为未来的推理优化研究提供了新的方向。

% ============================================================
\section{参考资料}
% ============================================================

\begin{enumerate}
  \item Leviathan, Y., et al. (2023). Fast Inference from Transformers via Speculative Decoding. ICML.
  \item Chen, C., et al. (2023). Accelerating Large Language Model Decoding with Speculative Sampling. arXiv.
  \item Kumar, T., et al. (2026). Speculative Speculative Decoding. arXiv:2603.03251.
\end{enumerate}

\end{document}